\backhead{FURTHER EXERCISES}{khred}
\begin{exlist}
\item A tilled layer has $\phi_0=0.58$, $\phi_c=0.46$ and $\tau_c=45$ days. Find the porosity after 15, 45 and 120 days, and the time to lose 90 per cent of the gain.\exans{0.546, 0.504, 0.468; 104 d}
\item A mulch saves 80 mm of evaporation in a season on the loam ($\psi_D=0.47$) and 50 mm on the sand ($\psi_D=0.71$). Find the recoverable water on each soil.\exans{42.4 mm; 14.5 mm}
\item A crop residue of 4 cm with tortuosity 2 lies over a wet soil. Find its resistance ($D_v=2.5\times10^{-5}$ m$^2$\,s$^{-1}$) and the fraction of evaporation remaining with $r_a=60$ and with $r_a=200$ s\,m$^{-1}$.\exans{3200 s\,m$^{-1}$; 0.018; 0.059}
\item A feedback loop has the links $+0.6$, $-0.5$ and $+0.4$. Find the gain, its sign and the amplification of a first perturbation.\exans{$-0.12$; negative; 0.89}
\item In the regime map, the drainage share of saved water is 0.3 at a rain scale of 0.6 and 0.8 at 1.4. Assuming a linear change, at which rain scale is half of the saved water drained?\exans{0.92}
\end{exlist}
\backhead{ESSAY QUESTIONS}{khred}
{\small\itshape Answer each question in about one page, with equations, a sketch and at least one number from this chapter.}\par\vspace{3pt}
\begin{exlist}
\item Explain how tillage and consolidation change the hydraulic properties of a soil, and how these changes affect irrigation.
\item Discuss the physics of evaporation under a mulch. Why does the saving depend on the soil and on the rain?
\item Explain the fate of the water saved by mulch and why a saving is not a gain until it is judged against a closed balance.
\item Describe the regime map and what its regions mean for a farmer who wants to save water.
\item Discuss the feedback loops between climate, management, soil water, crop and nitrogen shown in SWAFL. Which are positive and which negative?
\item Root structure changes the reach of the crop. Discuss how a deeper or denser root system changes the trigger and the dose.
\end{exlist}
\dykband{EXERCISES IN SWAFL}
{\small Open SWAFL (desktop or online) with the benchmark project. The chapter of the \href{https://github.com/halilunaziru73-creator/Coupled-Water-Fluxes-Pore-Scale-Dynamics-Precision-Irrigation-Mulching-Feedback-Loops/raw/main/Software/v1.1.0/SWAFL_Desktop/docs/manual_v1.1.0/SWAFL_Manual_v1.1.0.pdf}{SWAFL manual} that explains each task is given in brackets; record your values in a table and compare them with what is expected (\emph{Check}).}\par\vspace{3pt}
\begin{exlist}
\item {\color{khblue}[Manual, Chapter 8]} Open \emph{Tillage and structure}; compare no-till and tillage for the loam and record the drainage and the evaporation.\par\noindent{\small\itshape Check: tillage raises the infiltration and changes the drainage.}
\item {\color{khblue}[Manual, Chapter 9]} Open \emph{Mulch control}; vary the mulch strength from 0 to the maximum and record the saving of evaporation.\par\noindent{\small\itshape Check: the saving rises and then levels off.}
\item {\color{khblue}[Manual, Chapter 9]} Open \emph{Fate of saved water} and \emph{Regime map}; find the rain scale at which most of the saving drains.\par\noindent{\small\itshape Check: in wet climates most of the saving drains.}
\item {\color{khblue}[Manual, Chapter 10]} Open \emph{Feedback loops}; follow the seven effects and write the sentence under the picture in your own words.\par\noindent{\small\itshape Check: each effect has a sign and a size.}
\end{exlist}
